\documentclass[11pt]{article}

\usepackage[margin=1in]{geometry}
\usepackage{amsmath,amssymb}
\usepackage{booktabs}
\usepackage{hyperref}

\title{Localized Bump Response Threshold Release for HAOS-IIP}
\author{Tomislav Rupi\'c \\ Independent Researcher}
\date{April 2026}

\begin{document}
\maketitle

\begin{abstract}
This paper freezes release \texttt{53.2} of HAOS-IIP as a localized-bump response threshold result on the scalar carrier. The previous \texttt{52.5} smooth-inhomogeneity release left localized bump response open, with maximum refinement drift \texttt{0.5124}. Release \texttt{53.2} keeps the same scalar carrier and same localized bump profile, but excludes the earliest source-core transient layer using fit window \([0.010, 0.026]\). Under that bounded readout, all ten weak localized bump cases pass for amplitudes \(\eta = 0.01, 0.02, 0.03, 0.04, 0.05\), with maximum refinement drift \texttt{0.1095}, maximum median relative error \texttt{0.0943}, maximum p90 error \texttt{0.2984}, and maximum shell-\(\kappa\) CV \texttt{0.1248}. Stronger localized bumps remain open: the stress window \(\eta = 0.075, 0.10, 0.15\) reaches maximum drift \texttt{0.3578}. The correct conclusion is therefore thresholded: weak localized bump excitations close as a bounded metric-tracking transport response, while stronger localized excitations remain outside the same closure.
\end{abstract}

\tableofcontents

\section{Scope}

Release \texttt{53.2} is a narrow follow-on to the physics bridge and scalar-carrier transport stack.

It does not modify HAOS core, add new dynamics, claim particle ontology, or introduce stress-energy coupling.

It only asks:
\begin{quote}
does the localized bump response recover if the source-core transient layer is excluded from the constitutive fit?
\end{quote}

\section{Construction}

The test uses:
\begin{itemize}
\item the same scalar kernel-graph carrier;
\item the same localized bump shape with \(\sigma = 0.18\);
\item refinements \(n = 13, 15\);
\item delayed fit window \([0.010, 0.026]\);
\item the same current-closure thresholds used by the \texttt{52.5} inhomogeneity release.
\end{itemize}

Two amplitude regimes are separated:
\begin{itemize}
\item weak: \(\eta = 0.01, 0.02, 0.03, 0.04, 0.05\);
\item stress: \(\eta = 0.075, 0.10, 0.15\).
\end{itemize}

\section{Frozen Result}

\begin{table}[h!]
\centering
\begin{tabular}{@{}lccccc@{}}
\toprule
Regime & Cases pass & Max drift & Max median err. & Max p90 err. & Max shell CV \\
\midrule
weak & 10/10 & 0.1095 & 0.0943 & 0.2984 & 0.1248 \\
stress & 0/6 & 0.3578 & 0.2784 & 0.5287 & 0.3406 \\
\bottomrule
\end{tabular}
\caption{Weak localized bump response passes; stronger localized bumps remain open.}
\end{table}

The weakest passing weak case is:
\begin{itemize}
\item \texttt{weak\_bump|n=13|eta=0.050};
\item \(\kappa/D_{\mathrm{eff}} = 0.9699\);
\item median relative error \texttt{0.0943};
\item p90 relative error \texttt{0.2984};
\item shell-\(\kappa\) CV \texttt{0.1248};
\item \(|\mathrm{metric\ tracking\ corr}| = 0.9752\).
\end{itemize}

The hardest stress case is:
\begin{itemize}
\item \texttt{stress\_bump|n=13|eta=0.150};
\item median relative error \texttt{0.2784};
\item p90 relative error \texttt{0.5287};
\item shell-\(\kappa\) CV \texttt{0.3406}.
\end{itemize}

\section{Interpretation}

The old broad statement was:
\begin{quote}
localized bump response remains open.
\end{quote}

The updated statement is:
\begin{quote}
weak localized bump excitations close as a bounded metric-tracking transport response after the earliest source-core transient layer is excluded, while stronger localized bumps remain outside the same closure.
\end{quote}

This is closer to a proto-particle diagnostic than the previous result, but it is not a particle model. The result only establishes a weak localized-excitation proxy and an amplitude boundary on the scalar carrier.

\section{Artifacts}

\begin{itemize}
\item script: \path{numerics/simulations/scalar_kernel_graph_localized_bump_response.py};
\item operator note: \path{experiments/operators/Scalar_Kernel_Graph_Localized_Bump_Response_v1.md};
\item frozen result: \path{data/20260423_182539_scalar_kernel_graph_localized_bump_response.json};
\item latest result: \path{data/scalar_kernel_graph_localized_bump_response_latest.json};
\item summary plot: \path{plots/20260423_182539_scalar_kernel_graph_localized_bump_response_summary.png};
\item profile plot: \path{plots/20260423_182539_scalar_kernel_graph_localized_bump_response_profiles.png}.
\end{itemize}

\section{Next Boundary}

The next honest step is to refine the transition between \(\eta = 0.05\) and \(\eta = 0.075\), then test whether the weak localized response survives mild carrier disorder.

\section{Conclusion}

Release \texttt{53.2} turns the localized bump row from an undifferentiated failure into a thresholded result. Weak localized excitations now pass the scalar-carrier transport proxy; stronger localized excitations remain open. That is the correct stopping point: a recovered weak localized response, not a general localized-particle theory.

\end{document}
